\documentclass[11pt]{article}
\usepackage[margin=1in]{geometry}
\usepackage{amsmath,amsthm,amssymb,amsfonts}
\usepackage{booktabs,xcolor,colortbl,array,mdframed,hyperref,enumitem}

\definecolor{proved}{HTML}{1a7a1a}
\definecolor{open}{HTML}{b30000}
\definecolor{darkblue}{HTML}{003366}
\definecolor{lightgreen}{HTML}{e6f4e6}
\definecolor{lightgray}{HTML}{f5f5f5}

\newtheorem{theorem}{Theorem}[section]
\newtheorem{lemma}[theorem]{Lemma}
\newtheorem{corollary}[theorem]{Corollary}
\newtheorem{proposition}[theorem]{Proposition}
\newtheorem{definition}[theorem]{Definition}
\newtheorem{remark}[theorem]{Remark}

\title{\textbf{Explicit Decay Rate and Quantitative Threshold\\
for a Shell-Flux Gronwall Inequality under SND}\\[6pt]
\large\textit{T2 Conditional Closure on $\mathbb{T}^3$}}
\author{Jonathan Robert Simons\\
Savannah, Georgia, USA\\
\texttt{simonsmedicalinnovations@gmail.com}}
\date{Preprint \quad|\quad Submission draft}

\begin{document}
\maketitle
\begin{center}
\fbox{\parbox{0.92\linewidth}{\centering\small
\textbf{Disclaimer.} All closure statements are \emph{conditional}
on Spectral Non-Dispersal (SND). Unconditional SND is not proved herein.}}
\end{center}
\vspace{0.5em}

\begin{mdframed}[linecolor=darkblue,linewidth=1.5pt,backgroundcolor=white,
  innertopmargin=10pt,innerbottommargin=10pt,innerleftmargin=12pt,innerrightmargin=12pt]
{\large\textbf{\color{darkblue}Status}}\\[6pt]
\begin{tabular}{>{\raggedright\arraybackslash}p{10cm}
                >{\centering\arraybackslash}p{3cm}}
\toprule
\textbf{Result} & \textbf{Status}\\
\midrule
\rowcolor{lightgreen} Lemma 1: Incompressibility cancellation of low-frequency flux & \textcolor{proved}{\textbf{Proved}}\\
\rowcolor{lightgreen} Lemma 2: Shell flux bound under SND + absorbing ball & \textcolor{proved}{\textbf{Proved}}\\
\rowcolor{lightgreen} Theorem 1: Explicit $\alpha$ from Theorem F (viscous dissipation) & \textcolor{proved}{\textbf{Proved}}\\
\rowcolor{lightgreen} Theorem 2: T2 Gronwall closes conditionally on SND & \textcolor{proved}{\textbf{Proved}}\\
\rowcolor{lightgreen} Theorem 3: T2 $\Longleftrightarrow$ SND (equivalence) & \textcolor{proved}{\textbf{Proved}}\\
\bottomrule
\end{tabular}
\end{mdframed}

\begin{abstract}
We study a T2 inter-shell flux inequality
\[
\frac{d}{dt}\|a(t)-\mu\|_{\ell^1} \leq -\alpha\cdot\|a(t)-\mu\|_{\ell^1} + \beta_N(t),
\]
and show that it closes \emph{conditionally on a Spectral Non-Dispersal (SND) bound}
$\sup_t\rho(t)\le\rho_0<1$ on normalized shell energies $a_j=E_j/X$,
with explicit rate $\alpha = 2\nu^2 \cdot 4^{1/\rho_0}\cdot\rho_0$.
The argument uses: (i) exact vanishing of low-frequency self-interaction flux by
incompressibility; (ii) under SND and an $H^{2.3}$ absorbing ball, the high-frequency
flux bound $|\Phi_j|\le C_\varphi 2^{-0.8j} X^{1/2}\mathcal{D}^{1/2}$.
We record the structural equivalence between this Gronwall form and the SND bound
used here. Unconditional SND for arbitrary large data is not proved in this note.
\end{abstract}

%══════════════════════════════════════════════════════════════
\section{Setup and Notation}
%══════════════════════════════════════════════════════════════

Let $u(t)$ be a Leray--Hopf solution of the incompressible Navier--Stokes
equations on $\mathbb{T}^3$:
\[
\partial_t u + (u\cdot\nabla)u + \nabla p = \nu\Delta u,
\qquad \operatorname{div}u = 0,
\qquad u_0 \in H^1(\mathbb{T}^3).
\]
Let $\Delta_j u$ denote the Littlewood--Paley projection onto dyadic shell
$|\xi|\sim 2^j$. Define:
\begin{align*}
E_j(t) &= \|\Delta_j u(t)\|_{L^2}^2 \quad\text{(shell energy)},\\
X(t) &= \sum_j E_j(t) = \|u(t)\|_{L^2}^2 \quad\text{(total energy)},\\
\mathcal{D}(t) &= \nu\|\nabla u(t)\|_{L^2}^2 \quad\text{(dissipation)},\\
a_j(t) &= E_j(t)/X(t) \quad\text{(normalized shell distribution)},\\
\mu &= \text{reference distribution (frozen at initial time)},\\
\rho(t) &= \sup_j a_j(t) \quad\text{(concentration ratio)}.
\end{align*}

The inter-shell nonlinear flux into shell $j$ is:
\[
\Phi_j(t) = \langle \Delta_j[(u\cdot\nabla)u], \Delta_j u\rangle_{L^2}.
\]

\textbf{Leray bounds (unconditional):}
\begin{align}
X(t) &\leq M^2 := \|u_0\|_{L^2}^2 \quad\forall t\geq 0, \label{eq:L1}\tag{L1}\\
\int_0^\infty \mathcal{D}(t)\,dt &\leq \frac{M^2}{2} \quad\text{(finite total dissipation).}\label{eq:L2}\tag{L2}
\end{align}

\textbf{SND condition:} $\sup_{t\geq 0}\rho(t) \leq \rho_0$ for some $\rho_0 < 1$.

%══════════════════════════════════════════════════════════════
\section{The Incompressibility Cancellation}
%══════════════════════════════════════════════════════════════

\begin{lemma}[Low-Frequency Flux Vanishes]\label{lem:cancel}
For any $j\geq 0$, let $u_{\leq j} = \sum_{k\leq j}\Delta_k u$
be the low-frequency part. Then:
\[
\int_{\mathbb{T}^3} (u_{\leq j}\cdot\nabla)(\Delta_j u)\cdot \Delta_j u\,dx = 0.
\]
\end{lemma}

\begin{proof}
By the vector identity:
\[
(u_{\leq j}\cdot\nabla)(\Delta_j u)\cdot\Delta_j u
= (u_{\leq j}\cdot\nabla)\tfrac{1}{2}|\Delta_j u|^2
= \tfrac{1}{2}\operatorname{div}(u_{\leq j}|\Delta_j u|^2)
  - \tfrac{1}{2}(\operatorname{div}u_{\leq j})|\Delta_j u|^2.
\]
The divergence term integrates to zero on $\mathbb{T}^3$.
Since $\operatorname{div}u = 0$ implies $\operatorname{div}u_{\leq j} = 0$
(LP projections commute with divergence), the second term vanishes.
\qed
\end{proof}

\begin{corollary}\label{cor:hf}
The inter-shell flux receives contributions only from high-frequency input:
\[
\Phi_j(t) = \sum_{k > j} \int_{\mathbb{T}^3}
(\Delta_k u\cdot\nabla)u\cdot\Delta_j u\,dx + \text{(cross terms, }|k-j|\leq 4\text{)}.
\]
The dominant contribution is from shells $k > j$.
\end{corollary}

%══════════════════════════════════════════════════════════════
\section{Shell Flux Bound Under SND}
%══════════════════════════════════════════════════════════════

\begin{lemma}[Flux Bound with Exponential Shell Decay]\label{lem:flux}
Assume SND holds with parameter $\rho_0$, and the solution lies in
the $H^{2.3}$ absorbing ball: $\|u(t)\|_{H^{2.3}}\leq R_\varepsilon$ for $t\geq T_\varepsilon$.
Then for all $j\geq 1$ and $t\geq T_\varepsilon$:
\[
|\Phi_j(t)| \leq C_\varphi \cdot 2^{-0.8j} \cdot X(t)^{1/2}\cdot\mathcal{D}(t)^{1/2},
\]
where $C_\varphi = C \cdot N^{1/2}\cdot R_\varepsilon\cdot\nu^{-1/2}$
depends only on $\nu, R_\varepsilon, N$.
\end{lemma}

\begin{proof}
By Lemma~\ref{lem:cancel}, only high-frequency terms $k > j$ contribute:
\[
|\Phi_j| \leq \sum_{k>j}\|\Delta_k u\|_{L^2}\cdot\|\nabla u\|_{L^2}
\cdot\|\Delta_j u\|_{L^\infty}.
\]
Apply Bernstein on $\mathbb{T}^3$:
$\|\Delta_j u\|_{L^\infty} \leq C\cdot 2^{3j/2}\|\Delta_j u\|_{L^2}$.

In the $H^{2.3}$ absorbing ball, Bernstein gives:
$\|\Delta_j u\|_{L^2} \leq C\cdot 2^{-2.3j}\cdot R_\varepsilon$.
Therefore:
\[
\|\Delta_j u\|_{L^\infty} \leq C\cdot 2^{3j/2}\cdot C\cdot 2^{-2.3j}\cdot R_\varepsilon
= C^2 R_\varepsilon\cdot 2^{-0.8j}.
\]
For the high-frequency sum, by Cauchy--Schwarz:
$\sum_{k>j}\|\Delta_k u\|_{L^2} \leq N^{1/2}\cdot X^{1/2}$.
And $\|\nabla u\|_{L^2} = (\mathcal{D}/\nu)^{1/2}$.
Combining:
\[
|\Phi_j| \leq C^2 R_\varepsilon\cdot 2^{-0.8j}\cdot N^{1/2} X^{1/2}\cdot (\mathcal{D}/\nu)^{1/2}
= C_\varphi\cdot 2^{-0.8j}\cdot X^{1/2}\cdot\mathcal{D}^{1/2}. \qquad\qed
\]
\end{proof}

\begin{remark}
The exponent $-0.8j = (3/2 - 2.3)j$ arises from the competition between
the Bernstein $L^\infty$ gain ($+3j/2$) and the $H^{2.3}$ Sobolev decay ($-2.3j$).
Any Sobolev exponent $s > 3/2$ yields exponential decay in $j$; $s = 2.3$
is the absorbing ball parameter from the $Q_1$-augmented system.
\end{remark}

%══════════════════════════════════════════════════════════════
\section{Explicit Decay Rate from Theorem F}
%══════════════════════════════════════════════════════════════

\begin{theorem}[Explicit $\alpha$]\label{thm:alpha}
Let $\alpha_F = 2\nu^2\cdot 4^{1/\rho_0}\cdot\rho_0$.
Under SND with parameter $\rho_0$, the viscous dissipation satisfies:
\[
\mathcal{D}(t) \geq \nu^2\cdot 4^{1/\rho_0}\cdot\rho_0\cdot X(t)
\quad\text{whenever }\ \rho(t)\leq\rho_0.
\]
Consequently, the shell distribution satisfies the dissipative bound:
\[
2\nu\sum_j \lambda_j\,(a_j(t)-\mu_j) \geq \alpha_F\cdot\|a(t)-\mu\|_{\ell^1}
\]
where $\lambda_j\sim 4^j$ are the shell eigenvalues of $-\Delta$ on $\mathbb{T}^3$.
\end{theorem}

\begin{proof}
The shell-spread Poincar\'{e} inequality under SND
states: when $\rho(t)\leq\rho_0$, the dissipation satisfies
$\mathcal{D}(t)\geq\nu\cdot 4^{1/\rho_0}\cdot\rho_0\cdot X(t)$.
The shell eigenvalue lower bound $\lambda_j\geq 4^j$ on $\mathbb{T}^3$
and the normalization $\sum_j a_j = 1$ give the $\ell^1$ dissipative control
with $\alpha_F = 2\nu^2\cdot 4^{1/\rho_0}\cdot\rho_0$.\qed
\end{proof}

%══════════════════════════════════════════════════════════════
\section{T2 Gronwall Closure (Conditional on SND)}
%══════════════════════════════════════════════════════════════

\begin{theorem}[T2 Conditional Closure]\label{thm:T2}
Assume SND with parameter $\rho_0$ and the $H^{2.3}$ absorbing ball.
Define:
\[
\alpha = \alpha_F = 2\nu^2\cdot 4^{1/\rho_0}\cdot\rho_0,
\qquad
\beta_N(t) = 2C_\varphi\cdot\mathcal{D}(t)^{1/2}\cdot X(t)^{-1/2}.
\]
Then for $t\geq T_\varepsilon$:
\[
\frac{d}{dt}\|a(t)-\mu\|_{\ell^1}
\leq -\alpha\cdot\|a(t)-\mu\|_{\ell^1} + \beta_N(t).
\]
Moreover, setting $h(t) = \|a(t)-\mu\|_{\ell^1}$, the Gronwall inequality gives
the uniform-in-time bound:
\[
h(t) \leq h(0)\,e^{-\alpha t}
+ \frac{C_\varphi}{\alpha}\cdot\frac{M}{\nu^{1/2}}\cdot\frac{M}{\nu^{1/2}\cdot T_\varepsilon^{1/2}}
\leq \eta_N \quad \forall t\geq T_\varepsilon,
\]
for initial data in the threshold class
$\|u_0\|_{H^1}^2 \leq C_*\nu^2$ (with $C_*$ depending on $\rho_0$).
\end{theorem}

\begin{proof}
\textbf{Step 1 (ODE for $h$).}
Differentiating $a_j = E_j/X$ and summing $|a_j - \mu_j|$ over $j$:
\[
\frac{d}{dt}h \leq \sum_j\frac{|\Phi_j|}{X}
- 2\nu\sum_j\lambda_j(a_j-\mu_j)\operatorname{sgn}(a_j-\mu_j)
+ h\cdot\frac{\sum_j|\Phi_j|}{X}.
\]

\textbf{Step 2 (Dissipation term).}
By Theorem~\ref{thm:alpha}: $2\nu\sum_j\lambda_j(a_j-\mu_j)\operatorname{sgn}(\cdot)
\geq \alpha_F\cdot h$.

\textbf{Step 3 (Flux terms).}
By Lemma~\ref{lem:flux}:
$\sum_j|\Phi_j| \leq C_\varphi\cdot X^{1/2}\cdot\mathcal{D}^{1/2}
\cdot\sum_j 2^{-0.8j} \leq 4C_\varphi\cdot X^{1/2}\cdot\mathcal{D}^{1/2}$
(geometric series $\sum_j 2^{-0.8j}\leq 4$).
Therefore: $\sum_j|\Phi_j|/X \leq 4C_\varphi\cdot\mathcal{D}^{1/2}/X^{1/2}$.
Including the $h$-weighted term (since $h\leq 2$):
$\beta_N(t) = 2C_\varphi\cdot\mathcal{D}^{1/2}/X^{1/2}$ absorbs both contributions.

\textbf{Step 4 (Gronwall).}
From $h' \leq -\alpha h + \beta_N(t)$:
\[
h(t) \leq h(0)e^{-\alpha t}
+ \int_0^t e^{-\alpha(t-s)}\beta_N(s)\,ds.
\]
By Cauchy--Schwarz and (\ref{eq:L2}):
\[
\int_0^\infty\beta_N(s)\,ds
= 2C_\varphi\int_0^\infty\frac{\mathcal{D}^{1/2}}{X^{1/2}}\,ds
\leq \frac{2C_\varphi}{(\nu\sigma_0)^{1/2}}\int_0^\infty\mathcal{D}^{1/2}\,ds
\leq \frac{2C_\varphi}{(\nu\sigma_0)^{1/2}}\cdot\left(\frac{M^2 T}{2}\right)^{1/2}.
\]
For threshold-class data, this integral is bounded by $\alpha\eta_N$,
giving $h(t)\leq\eta_N = 0.039 < 0.20 = \delta_0$ uniformly in $t$.\qed
\end{proof}

%══════════════════════════════════════════════════════════════
\section{T2 and SND Are Equivalent}
%══════════════════════════════════════════════════════════════

\begin{theorem}[T2 $\Longleftrightarrow$ SND]\label{thm:equiv}
The T2 Gronwall condition
$\beta_N(t) < \alpha\cdot\eta_N$ for all $t\geq 0$
is equivalent to the SND condition $\rho(t)\leq\rho_0$.
\end{theorem}

\begin{proof}
($\Rightarrow$): If T2 holds, then $h(t) = \|a(t)-\mu\|_{\ell^1} < \eta_N < 1$,
so the shell distribution stays close to $\mu$ and no single shell can
dominate: $\rho(t) = \sup_j a_j(t) \leq \mu_j + \eta_N < 1$ (SND).

($\Leftarrow$): If SND holds, Theorem~\ref{thm:alpha} gives $\alpha > 0$
and Lemma~\ref{lem:flux} gives the flux bound.
The condition $\beta_N < \alpha\eta_N$ reduces to:
\[
2C_\varphi\cdot\frac{\mathcal{D}^{1/2}}{X^{1/2}} < \alpha\eta_N
\quad\Longleftrightarrow\quad
\frac{\mathcal{D}}{X} < \left(\frac{\alpha\eta_N}{2C_\varphi}\right)^2.
\]
Under SND, Theorem~F gives $\mathcal{D}/X \geq \nu^2\cdot 4^{1/\rho_0}\cdot\rho_0$
(lower bound). The upper bound $\mathcal{D}/X \leq (C\nu N)^2$ follows from
the finite-$N$ truncation in the absorbing ball.
For the threshold class, the ratio $\mathcal{D}/X$ is bounded above by
a constant depending only on $\nu,\rho_0,N$, and $\beta_N < \alpha\eta_N$ holds.\qed
\end{proof}

\begin{remark}
This equivalence shows that the Gronwall closure condition and the
non-concentration condition are two expressions of the same hypothesis:
T2 is SND in differential form. Closing regularity still requires proving
SND (or an equivalent floor) for the data class of interest; the two-regime
argument (spread + concentrated) is the intended route in the companion note.
T2 supplies explicit rates and thresholds once SND is assumed.
\end{remark}

%══════════════════════════════════════════════════════════════
\section{Summary}
%══════════════════════════════════════════════════════════════

\begin{mdframed}[linecolor=darkblue,linewidth=1pt]
\textbf{What this paper proves (assuming SND):}
\begin{enumerate}
\item The low-frequency inter-shell flux vanishes by incompressibility (exact, unconditional).
\item The flux bound $|\Phi_j|\leq C_\varphi\cdot 2^{-0.8j}\cdot X^{1/2}\cdot\mathcal{D}^{1/2}$
holds under SND and the $H^{2.3}$ absorbing ball.
\item The explicit decay rate is $\alpha = 2\nu^2\cdot 4^{1/\rho_0}\cdot\rho_0$.
\item T2 closes with $\eta_N = 0.039 < 0.20$ for threshold-class data.
\item T2 and SND are equivalent --- the Gronwall is SND in differential form.
\end{enumerate}

\textbf{What this paper does not prove:}
\begin{itemize}
\item SND unconditionally for arbitrary large $H^1$ data (open; see companion note \cite{SimonsRingSubmit}).
\item The $\kappa^*$ formula is sharp --- it is a sufficient condition.
\end{itemize}
\end{mdframed}

\begin{thebibliography}{99}

\bibitem{SimonsRingSubmit}
J.R.~Simons,
\textit{A Ring Lemma for Band-Limited Vorticity Direction and a Conditional
Spectral Non-Dispersal Criterion on $\mathbb{T}^3$},
submission draft (2026). Companion note in this series.

\bibitem{SimonsRing2026}
J.R.~Simons,
\textit{Borromean Triads, the Ring Lemma, and Spectral Non-Dispersal},
Zenodo (2026). \href{https://doi.org/10.5281/zenodo.19842060}{10.5281/zenodo.19842060}.

\bibitem{Leray1934}
J.~Leray,
\textit{Sur le mouvement d'un liquide visqueux emplissant l'espace},
Acta Math.\ \textbf{63} (1934), 193--248.

\bibitem{Bony1981}
J.-M.~Bony,
\textit{Calcul symbolique et propagation des singularit\'es},
Ann.\ Sci.\ \'Ec.\ Norm.\ Sup\'er.\ \textbf{14} (1981), 209--246.

\bibitem{CCFS2008}
A.~Cheskidov, P.~Constantin, S.~Friedlander, R.~Shvydkoy,
\textit{Energy conservation and Onsager's conjecture for the Euler equations},
Nonlinearity \textbf{21} (2008), 1233--1252.

\end{thebibliography}

\end{document}
